\documentclass[11pt]{article}

\usepackage[T1]{fontenc}
\usepackage[a4paper,margin=1in]{geometry}
\usepackage{amsmath,amssymb,bm}
\usepackage{booktabs,tabularx}
\usepackage[hidelinks]{hyperref}
\makeatletter
\g@addto@macro\UrlBreaks{\do\-\do\_}
\makeatother
\Urlmuskip=0mu plus 1mu

\setlength{\parindent}{0pt}
\setlength{\parskip}{0.65em}
\newcommand{\method}[1]{\texttt{#1}}

\title{\vspace{-1em}TS-IFA: Current Model and Experimental Design\vspace{-0.5em}}
\author{}
\date{}

\begin{document}
\maketitle

\section{Overall goal and mathematical setting}

The goal of TS-IFA is to adapt a frozen pretrained forecaster at inference time
using retrieved historical evidence, without updating the backbone parameters.
The repository compares this learned adapter to simpler output-space adapters,
learned gates, and target-aware oracles.  All methods take the same basic
objects: a query history, the frozen model forecast, and $K$ retrieved complete
historical windows.

For a query lookback $x\in\mathbb{R}^{L}$ and target
$y\in\mathbb{R}^{H}$, a frozen forecaster $f$ produces the vanilla forecast
\[
  \hat y=f(x).
\]
The implementation may also supply dataset covariates to $f$; they are omitted
from the notation.  Retrieval returns $K$ complete historical examples
\[
  \mathcal N(x)=\{(x_k,y_k,d_k)\}_{k=1}^{K},
  \qquad x_k\in\mathbb{R}^{L},\quad y_k\in\mathbb{R}^{H},
\]
ordered by distance $d_k$.  When the backbone exposes a context-conditioned
mode, the scaled trajectories $[x_k,y_k]$ can be passed as context channels to
obtain
\[
  \hat y^c=f\!\left(x;\{[x_k,y_k]\}_{k=1}^{K}\right).
\]
Each neighbor is also forecast without retrieved context,
$\hat y_k=f(x_k)$, and its realized residual is
\[
  e_k=y_k-\hat y_k.
\]
Thus extraction computes one vanilla query forecast, one context-conditioned
query forecast, and the $K$ neighbor forecasts needed by residual methods.

Retrieval is represented by a feature map $\psi$ and a distance $\rho$:
\[
\mathcal N(x)=\operatorname{KNN}_K
\left(\psi(x),\{\psi(x_i)\}_i;\rho\right),\qquad
\psi(x)=
\begin{cases}
x, & \text{raw},\\[2mm]
(x-\mu_x)/(\sigma_x+\varepsilon), & \text{instance-normalized},\\[2mm]
g(x), & \text{encoder ablation},
\end{cases}
\]
where $g(x)$ is a backbone representation when encoder-space retrieval is
requested.  Retrieval is shape-based, but the adapters operate on the query's
level and scale.  Let $(\mu_x,\sigma_x)$ and $(\mu_k,\sigma_k)$ be statistics
of the query and neighbor lookbacks.  Before a neighbor value $v_k$ (lookback,
horizon, or forecast) is used, it is transformed as
\[
  \mathcal S_{k\rightarrow x}(v_k)
  =\frac{v_k-\mu_k}{\sigma_k}\sigma_x+\mu_x.
\]
Residuals contain no additive level and receive only the scale transform,
\[
  \mathcal S^{\mathrm{res}}_{k\rightarrow x}(e_k)
  =\frac{e_k}{\sigma_k}\sigma_x.
\]
All neighbor quantities below refer to these query-scale versions.

\section{Output-space adapters and baselines}

For distance weighting, the code first normalizes distances within each query,
\[
  \bar d_k=\frac{d_k-\min_j d_j}{\max(\operatorname{std}_j d_j,\varepsilon)},
  \qquad
  w_k=\frac{\exp(-\bar d_k)}{\sum_j\exp(-\bar d_j)}.
\]
The following non-trainable baselines are evaluated:
\begin{align*}
\text{context conditioned:}\quad &\hat y^c,\\
\text{unweighted neighbor mean:}\quad &K^{-1}\sum_k y_k,\\
\text{weighted neighbor mean:}\quad &\sum_k w_k y_k,\\
\text{weighted residual correction:}\quad
  &\hat y+\sum_k w_k e_k.
\end{align*}

Three trainable mixtures are fitted on a calibration split by ridge regression
on the residual target $y-\hat y$:
\begin{align*}
\text{convex weighted mix:}\quad
  \hat y^{(0)}
  &=(1-\lambda)\hat y+\lambda\sum_k w_k y_k,
  &&\lambda\in[0,1],\\
\text{shared learned mix:}\quad
  \hat y^{(1)}_h
  &=\hat y_h+a_0\hat y_h+\sum_k a_k y_{k,h},\\
\text{horizon-specific full mix:}\quad
  \hat y^{(2)}_h
  &=\hat y_h+a_{0,h}\hat y_h+
    \sum_k a_{k,h}y_{k,h}+b_h\hat y^c_h.
\end{align*}
The regressions have no intercept, so zero coefficients recover the vanilla
forecast.  Features are RMS-standardized without centering, and the normal
equations are averaged over observations before adding the $\ell_2$
coefficient.  This keeps regularization comparable across dataset units and
payload sizes.

\section{Learned gates and target-aware oracles}
\label{sec:gates}

The gate family selects between $\hat y$ and $\hat y^c$.  The implemented
learned gates use classifier and signed-loss-regressor objectives, each with a
scalar and a horizon-wise output.

For a target $y$, define the improvement of context over vanilla at horizon $h$
as
\[
  \Delta_h=(y_h-\hat y_h)^2-(y_h-\hat y^c_h)^2.
\]
The scalar target is $H^{-1}\sum_h\Delta_h$; the horizon target is the vector
$(\Delta_h)_{h=1}^{H}$.  Classifiers predict whether the relevant improvement
is positive and select context at probability greater than $0.5$.  Regressors
predict the signed improvement and select context when it is positive.
Scalar gates make one decision for the complete forecast, whereas horizon-wise
gates make one decision per horizon step.

All gates use information available at inference.  Scalar models receive the
mean and standard deviation of $\hat y^c-\hat y$; horizon models receive its
full $H$-vector.  Both receive 13 common features: the mean weighted-neighbor
minus vanilla forecast, mean weighted residual, query mean and standard
deviation, four summaries of neighbor-lookback means and standard deviations,
same-series neighbor ratio, mean neighbor age, standard deviation and maximum
of the neighbor weights, and mean retrieval distance.  Classifier gates use
balanced class weights.

Two target-aware upper bounds are also evaluated:
\begin{align*}
\text{scalar oracle:}\quad
  &\text{choose the lower-MSE forecast between $\hat y$ and $\hat y^c$ for each sample},\\
\text{horizon oracle:}\quad
  &\text{choose the lower-squared-error forecast independently for each horizon step}.
\end{align*}
These methods use the target to make their decision.  They are therefore
oracles, not deployable gates, and are reported only in the dedicated gate
table.

\section{TS-IFA model}

TS-IFA is a lightweight output-space adapter.  Its inputs are $x$, $\hat y$,
$\hat y^c$, and the query-scale neighbor tensors $X_c$, $Y_c$, $\hat Y_c$, and
$E_c=Y_c-\hat Y_c$.  The frozen Chronos model is not part of TS-IFA training.

The implementation uses a projected multi-head cross-attention block.  Each
block linearly projects and layer-normalizes queries, keys, and values; applies
multi-head attention; adds query and attention residuals followed by layer
normalization; applies a four-times-wider feed-forward sublayer with another
residual and normalization; and finally projects to the requested output size.
The forward pass returns attention weights; evaluation artifacts omit them.

Let $m=[x,\hat y]$ and $m_k=[x_k,\hat y_k]$.  The residual branch is
\[
  z_r=\operatorname{CrossAttn}_r(m,\{m_k\},\{e_k\}),\qquad
  y_r=\hat y+\operatorname{MLP}_r(z_r).
\]
It learns which historical forecast errors are relevant to the query.  The
memory branch is
\[
  z_m=\operatorname{CrossAttn}_m(x,\{x_k\},\{y_k\}),\qquad
  y_m=\operatorname{MLP}_m([\hat y,z_m]).
\]
It learns directly from realized neighbor horizons.

The four candidate forecasts are
\[
  \widetilde Y=[\hat y,\hat y^c,y_r,y_m]\in\mathbb{R}^{4\times H}.
\]
The query and each candidate are embedded independently,
\[
  z_x=\operatorname{MLP}_x([x,\hat y]),\qquad
  z_{y,j}=\operatorname{MLP}_y(\widetilde Y_j),
\]
and a third cross-attention block produces a learned adaptation proposal,
\[
  \tilde y_{\mathrm{mix}}
  =\operatorname{CrossAttn}_{\mathrm{mix}}(z_x,\{z_{y,j}\}_{j=1}^{4},\widetilde Y).
\]
This proposal is wrapped in an explicit identity-preserving residual gate.  A
horizon-wise gate is computed from the query, the vanilla forecast, and the
mixture proposal:
\[
  \alpha
  =\sigma\!\left(\operatorname{MLP}_{\alpha}([x,\hat y,\tilde y_{\mathrm{mix}}])\right)
  \in (0,1)^H,
\]
and the final prediction is
\[
  \hat y_{\mathrm{TS\mbox{-}IFA}}
  =\hat y+\alpha\odot(\tilde y_{\mathrm{mix}}-\hat y).
\]
Thus $\alpha=0$ gives the exact vanilla forecast, while $\alpha=1$ gives the
full learned mixture proposal.  The gate MLP is initialized with zero final
weights and bias $-6$, so the initial gate value is
$\sigma(-6)\simeq 2.5\times 10^{-3}$ and training starts very close to the
identity map.

All three attention blocks default to four heads with per-head dimension 32.
The residual, memory, mixture, and final gate MLP hidden widths are 128, the
mixture key dimension is 64, the default mixture gate bias is $-6$, and dropout
is zero.  The implementation records exact total and trainable parameter counts
in both the checkpoint and
\method{config.json}; this is preferable to a simplified count that omits the
projection, normalization, and feed-forward parameters inside each attention
block.

\section{TS-IFA preprocessing and training}

After the neighbor-to-query scale transfer, TS-IFA can apply a second instance
normalization using the query statistics.  Query values, forecasts, and
neighbor values receive both location and scale normalization; neighbor
residuals receive scale normalization only.

The optimized loss is
\[
  \mathcal L
  =\underbrace{\operatorname{MSE}(\hat y_{\mathrm{TS\mbox{-}IFA}},y)}_{\mathcal L_{\mathrm{pred}}}
  +\beta\underbrace{\operatorname{MSE}(\hat y_{\mathrm{TS\mbox{-}IFA}},\hat y)}_{\mathcal L_{\mathrm{reg}}}
  +\gamma\underbrace{\operatorname{MSE}(\delta_r,y-\hat y)}_{\mathcal L_r},
\]
where $\delta_r=\operatorname{MLP}_r(z_r)$.  With instance normalization these
are normalized MSE terms.  Without it, each error is divided by the query
lookback standard deviation before averaging.

\section{Experimental protocol and retrieval implementation}

This section gives the current experimental choices.  It is intentionally
separated from the mathematical description above: the formulas define the
adapters, while the details below define the present cluster runs.

\subsection{Chronological protocol}

Every dataset is split chronologically with ratios $(0.30,0.35,0.15,0.20)$:

\begin{center}
\begin{tabularx}{\textwidth}{@{}c l X@{}}
\toprule
Period & Role & Consumers \\
\midrule
$T_0$ & Initial datastore & Fixed-retrieval reference and online-capacity reference \\
$T_1$ & Baseline training & Learned mixtures; first part of TS-IFA training \\
$T_2$ & Gate training & Learned context gates; second part of TS-IFA training \\
$T_3$ & Final evaluation & Direct Chronos, held-out baselines, gates, oracles, and TS-IFA \\
\bottomrule
\end{tabularx}
\end{center}

The saved payload names are \method{train}, \method{oracle}, and
\method{eval}, corresponding to $T_1$, $T_2$, and $T_3$.  The name
\method{oracle} is historical: this payload is the training data for the
learned gates.  Only the explicitly named target-aware methods in
Section~\ref{sec:gates} are oracles.

Mixtures are fitted only on $T_1$.  Gates are fitted only on $T_2$.  TS-IFA is
trained on the concatenation $T_1\cup T_2$, without a validation split or early
stopping.  All reported final metrics are evaluated on untouched $T_3$ data,
except for the clearly suffixed \method{\_eval\_fit} diagnostic bounds.  The
\method{\_eval\_fit} mixture results are fitted directly on $T_3$ targets and
are optimistic in-sample bounds, not held-out forecasting methods; they are
excluded from the main table and from best-value bolding.

Temporal sampling is split into four explicit stride controls.  The datastore
stride controls the phase-aligned retrieval history; when period alignment is
enabled it must be a multiple of the alignment period, otherwise later
datastore dates would drift away from the query phase.  The $T_1$, $T_2$, and
$T_3$ query payloads are controlled separately by the train, oracle, and eval
strides, respectively.

\subsection{Retrieval implementation}

Candidate windows are aligned to the query phase (period 24 by default), and a
candidate is eligible only when its complete lookback and horizon precede the
query lookback.  Windows from every target series are pooled in the datastore,
so a neighbor may come from the query series or another series.

Exact $K$-nearest-neighbor search is implemented in raw lookback space,
instance-normalized lookback space, and encoder space.  Encoder features can
optionally be pooled.  Euclidean, cosine, and Pearson distances are available.
Search is exact, with query chunking used only to bound memory.

The current cluster experiments use Chronos-2, ten neighbors, exact Euclidean
nearest-neighbor search on independently instance-normalized time-domain
lookbacks, and online retrieval:
\[
  K=10,\qquad d_k=\lVert\psi(x)-\psi(x_k)\rVert_2,
  \qquad \psi=\text{instance-normalized lookback}.
\]
Nearest-neighbor search in the Chronos encoder space is therefore
\emph{not} the default.  Raw-lookback and encoder-space retrieval remain
implemented as ablations.

In the default \method{online} mode, each query uses its most recent eligible
history.  The number of aligned dates is capped by the corresponding $T_0$
capacity unless \method{--full-online-history} is set; the current jobs also cap
the store at 30,000 windows.  Consequently, later queries may use observations
from $T_1$ or $T_2$ when these observations are already strictly in their past.
In the \method{fixed} ablation, every $T_1$, $T_2$, and $T_3$ query retrieves
only from $T_0$.  Date and window caps, explicit history bounds, the minimum
store size, and period alignment are independently configurable.

\subsection{Current TS-IFA training run}

The current cluster job trains for 200 epochs on $T_1\cup T_2$ with AdamW,
batch size 256, learning rate $10^{-3}$, weight decay $10^{-4}$, gradient-norm
clipping at 1, $\beta=10^{-2}$, and $\gamma=0$.  Hence residual-branch
supervision is implemented but disabled in the current experiment.  There is
no validation evaluation during training; $T_3$ is evaluated once after the
final epoch.

\section{Current experiment configuration}

\begin{center}
\begin{tabularx}{\textwidth}{@{}l X@{}}
\toprule
Component & Current cluster value \\
\midrule
Datasets & Electricity; Solar aggregated by hourly sum \\
Window settings & $(168,24)$ and $(672,168)$ \\
Backbone & Chronos-2, frozen, local weights, future neighbor context enabled \\
Splits & $T_0/T_1/T_2/T_3=0.30/0.35/0.15/0.20$ \\
Retrieval & online, phase period 24, exact Euclidean in instance-normalized lookback space \\
Neighbors and capacity & $K=10$, at most 30,000 datastore windows \\
Sampling & datastore stride 24; $T_1$ query stride 24; $T_2$ query stride 24; $T_3$ query stride 128 \\
Baseline ridge & $\ell_2=10^{-3}$; mixtures fitted on $T_1$ \\
Gates & CatBoost, 300 iterations, learning rate 0.03, depth 4; fitted on $T_2$ \\
TS-IFA & 200 epochs, batch 256, AdamW, learning rate $10^{-3}$, weight decay $10^{-4}$; final gate initialized near identity \\
TS-IFA loss & instance normalization, $\beta=10^{-2}$, $\gamma=0$ \\
Seed & 1 \\
Final evaluation & $T_3$ only, using nMSE in generated result tables \\
\bottomrule
\end{tabularx}
\end{center}

\section{Code organization and experiment entry points}

The repository is organized as follows:
\begin{itemize}
  \item \path{ts_ifa/data/load_dataset_model.py}: CSV loading, chronological
  split bounds, model construction, device handling, and run directories.
  \item \path{ts_ifa/data/neighbors.py}: aligned datastore dates, raw/instance/
  encoder representations, and chunked exact nearest-neighbor search.
  \item \path{ts_ifa/data/scaling.py}: neighbor-to-query scale transfer used by
  baselines and TS-IFA.
  \item \path{ts_ifa/models/chronos_model.py}: frozen Chronos-2 wrapper,
  optional past/future context channels, and optional encoder representations.
  \item \path{ts_ifa/models/ts_ifa.py}: attention block, \method{TSIFAConfig},
  residual and memory branches, final mixture proposal, and identity-initialized
  residual gate.
  \item \path{ts_ifa/experiments/experiment_univariate.py}: direct backbone
  evaluation without retrieval.
  \item \path{ts_ifa/experiments/extraction.py}: the expensive shared stage;
  it retrieves neighbors and writes $T_1/T_2/T_3$ prediction and feature payloads.
  \item \path{ts_ifa/experiments/evaluate_baselines.py}: scale transfer,
  non-trainable baselines, ridge mixtures, CatBoost gates, and true context oracles.
  \item \path{ts_ifa/experiments/train_ts_ifa.py}: payload dataset, optional
  instance normalization, TS-IFA optimization, final evaluation, and artifacts.
  \item \path{ts_ifa/results_table.py}: aggregation of direct, baseline, gate,
  and TS-IFA metrics into LaTeX tables.
  \item \path{ts_ifa/visu/retrieval_dashboard.ipynb}: artifact-only inspection
  of examples, errors, gate scores, and decisions; it does not retrain models.
\end{itemize}

The cluster workflow is deliberately staged.  First, the
\path{extract_payloads.slurm} launcher under \path{ts_ifa/slurm/extraction/}
evaluates direct Chronos forecasts and creates the shared retrieval payloads.  Once it succeeds, the
baseline, gate, and TS-IFA jobs can run concurrently from
\path{ts_ifa/slurm/baselines/}, \path{ts_ifa/slurm/gates/}, and
\path{ts_ifa/slurm/training/}.  Finally,
the \path{build_results_table.slurm} launcher under
\path{ts_ifa/slurm/results/} creates the main,
baseline, gate, and TS-IFA tables.  Downstream stages read the same extraction
directory but write to separate subdirectories, so their artifacts do not
overwrite one another.

\subsection{Important configuration and ablation switches}

The literal configuration blocks at the top of the SLURM scripts are the
source of truth for cluster experiments.  The most important ablation axes are:

\begin{itemize}
  \item \textbf{Retrieval space:} \path{DISTANCE_SPACES} and
  \path{--distance-space} select \method{instance} (default), \method{raw}, or
  \method{encoder}; \path{--pool-representation} enables encoder pooling.
  \item \textbf{Distance metric:} \path{--distance-metric} selects
  \method{euclidean} (current), \method{cosine}, or \method{pearson}.
  \item \textbf{Neighbor count:} \path{NEIGHBORS} and \path{--neighbors}
  control payload shape and therefore require a new extraction run.
  \item \textbf{History policy and capacity:} \path{RETRIEVAL_MODE} and
  \path{--retrieval-mode} select \method{online} (current) or fixed-$T_0$
  \method{fixed}.  Date/window caps use \path{--max-store-dates} and
  \path{--max-store-windows}; \path{--full-online-history} removes the implicit
  $T_0$ date cap.
  \item \textbf{Temporal sampling:} \path{--period} controls phase alignment;
  \path{--datastore-stride} controls aligned retrieval history density;
  \path{--train-stride}, \path{--oracle-stride}, and \path{--eval-stride}
  control $T_1$, $T_2$, and $T_3$ query density.
  \item \textbf{Context residuals:} \path{--compute-ec} adds
  neighbor-context residuals to the payload.  It is disabled because current
  baselines and TS-IFA use univariate neighbor residuals $e_k$.
  \item \textbf{Gate capacity:} \path{GATE_ITERATIONS},
  \path{GATE_LEARNING_RATE}, and \path{GATE_DEPTH} control CatBoost iterations,
  learning rate, and depth.
  \item \textbf{TS-IFA architecture:} the training CLI exposes
  \path{--residual-heads}, \path{--memory-heads}, \path{--mixture-heads}, three
  attention-dimension options, hidden widths, mixture key dimension,
  \path{--mixture-gate-init}, and dropout.
  \item \textbf{TS-IFA objective:} \path{BETA}, \path{GAMMA}, and
  \path{--normalization} control conservativeness, explicit residual
  supervision, and instance-normalization ablations.
\end{itemize}

Retrieval-space, metric, neighbor-count, history-policy, capacity, or sampling
changes alter the payload and must be made in the extraction stage first.  The
same run name and settings must then be used by every downstream job.  Gate-only
and TS-IFA-only hyperparameters consume existing payloads and can be varied
without rerunning Chronos.

\section{Artifacts and checks}

Extraction writes one \path{*_prediction_payload.pt} file per extracted split,
plus feature payloads and a retrieval example.  Baselines and gates
write JSON/CSV metrics, fitted artifacts, and compact visualization payloads.
TS-IFA writes the checkpoint, configuration, training history, final metrics,
evaluation predictions, and training curve under its run subdirectory.

The lightweight local checks most relevant to this document are
\path{tests/smoke/check_baseline_oracles.py} and
\path{tests/smoke/check_ts_ifa_training.py}.  Full extraction and training are
cluster workloads and should not be run locally.

\end{document}
